\documentclass[12pt]{article}
\usepackage[utf8]{inputenc}
\usepackage[margin=1in]{geometry}
\usepackage{amsmath}\usepackage{amssymb}\usepackage{multicol}\usepackage{enumitem}
\usepackage{pgfplots}\pgfplotsset{compat=1.17}
\setlength{\parindent}{0pt}
\begin{document}

\fbox{\parbox{0.97\linewidth}{\small\textbf{Final answers} (record here --- graded on the answer; show your work):\\[6pt]1.~\underline{\hspace{1.5cm}}\quad 2.~\underline{\hspace{1.5cm}}\quad 3.~\underline{\hspace{1.5cm}}\quad 4.~\underline{\hspace{1.5cm}}\quad 5.~\underline{\hspace{1.5cm}}\quad 6.~\underline{\hspace{1.5cm}}\quad 7.~\underline{\hspace{1.5cm}}\quad 8.~\underline{\hspace{1.5cm}}\quad 9.~\underline{\hspace{1.5cm}}\quad 10.~\underline{\hspace{1.5cm}}\quad 11.~\underline{\hspace{1.5cm}}}}
\vspace{0.35cm}

\begin{center}{\Large \textbf{Even, Odd, or Neither}}\end{center}\vspace{0.2cm}
% Determine even, odd, or neither
\textbf{Example (Determine even, odd, or neither):}

$f(x)=x^{4}-x^{2}$: $f(-x)=(-x)^4-(-x)^2=x^4-x^2=f(x)$, so it is \textbf{even}.

\vspace{0.2cm}\hrule\vspace{0.35cm}
\textbf{Directions: Compute $f(-x)$ and simplify, then classify as even, odd, or neither.}

\begin{multicols}{2}
\begin{enumerate}[label=\textbf{\arabic*.}, start=1, itemsep=2cm]
  \item $\displaystyle f(x)=x^{2}+1$
  \item $\displaystyle f(x)=x^{3}-x$
  \item $\displaystyle f(x)=x^{3}+2$
  \item $\displaystyle f(x)=|x|$
  \item $\displaystyle f(x)=2x$
  \item $\displaystyle f(x)=x^{4}-3x^{2}$
  \item $\displaystyle f(x)=x^{5}$
  \item $\displaystyle f(x)=x^{2}+x$
  \item $\displaystyle f(x)=\dfrac{1}{x}$
  \item \textbf{(challenge)}\ $\displaystyle f(x)=x^{3}+x^{2}$
\end{enumerate}
\end{multicols}

\vspace{0.2cm}\hrule\vspace{0.3cm}
\textbf{Challenge} (same sheet):\\[4pt]
\begin{enumerate}[label=\textbf{\arabic*.}, start=11, itemsep=2.2cm]
  \item \textbf{(challenge)}\ Explain why any function made ONLY of even-power terms (like $x^4-3x^2+5$) must always be even, using the substitution rule.
\end{enumerate}

\newpage

\begin{center}{\Large \textbf{Answer Key --- fully worked}}\end{center}
\vspace{0.25cm}
\textbf{Procedure:} Replace every $x$ in the function with $-x$.; Simplify the result completely.; If $f(-x)=f(x)$, the function is even.; If $f(-x)=-f(x)$, the function is odd.; If neither matches, the function is neither..\\[8pt]
\begin{enumerate}[label=\textbf{\arabic*.},itemsep=7pt]
  \item $f(-x)=x^{2}+1=f(x)$: even
  \item $f(-x)=-x^{3}+x=-(x^{3}-x)=-f(x)$: odd
  \item $f(-x)=-x^{3}+2$, which is neither $f(x)$ nor $-f(x)$: neither
  \item $f(-x)=|-x|=|x|=f(x)$: even
  \item $f(-x)=-2x=-f(x)$: odd
  \item $f(-x)=x^{4}-3x^{2}=f(x)$: even (all even-power terms)
  \item $f(-x)=-x^{5}=-f(x)$: odd
  \item $f(-x)=x^{2}-x$, which is neither $f(x)$ nor $-f(x)$: neither
  \item $f(-x)=\dfrac{1}{-x}=-\dfrac{1}{x}=-f(x)$: odd
  \item $f(-x)=-x^{3}+x^{2}$: neither --- the even $x^2$ term stays positive while the odd $x^3$ term flips sign, so it can't equal $f(x)$ or $-f(x)$
  \item When you substitute $-x$ into any even-power term, the negative disappears: $(-x)^2=x^2$, $(-x)^4=x^4$, and in general $(-x)^{2n}=x^{2n}$ because a negative base raised to an even power is positive. A constant term (like $+5$) is also unchanged. So if EVERY term has an even power (or is a constant), then $f(-x)$ equals the original $f(x)$ term-for-term. For example, $f(x)=x^4-3x^2+5$ gives $f(-x)=(-x)^4-3(-x)^2+5=x^4-3x^2+5=f(x)$. Since $f(-x)=f(x)$, the function must be even.
\end{enumerate}

\end{document}